\documentclass{article}
\usepackage{amsmath, amssymb, amsthm}
\usepackage{hyperref}
\usepackage{geometry}
\geometry{margin=1in}

\title{Erd\H{o}s Problem \#307}
\date{Last updated: 2025-09-09}
\author{Source: erdosproblems.com}

\begin{document}
\maketitle

\noindent\textbf{Status:} VERIFIABLE \quad \textbf{No prize}

\noindent\textbf{Tags:} number theory, unit fractions

\medskip

\noindent\textbf{Problem Statement:}

Are there two finite sets of primes $P,Q$ such that\[1=\left(\sum_{p\in P}\frac{1}{p}\right)\left(\sum_{q\in Q}\frac{1}{q}\right)?\]




Asked by Barbeau \cite{Ba76}. Can this be done if we drop the requirement that all $p\in P$ are prime and just ask for them to be relatively coprime, and similarly for $Q$? Cambie has found several examples when this weakened version is true. For example,\[1=\left(1+\frac{1}{5}\right)\left(\frac{1}{2}+\frac{1}{3}\right)\]and\[1=\left(1+\frac{1}{41}\right)\left(\frac{1}{2}+\frac{1}{3}+\frac{1}{7}\right).\]There are no examples known of the weakened coprime version if we insist that $1\not\in P\cup Q$. It is easy to see that, if $P$ and $Q$ are sets of primes, then $P$ and $Q$ are disjoint, and $\sum_{p\in P\cup Q}\frac{1}{p}\geq 2$, whence $\lvert P\cup Q\rvert \geq 60$.




References

[Ba76] Barbeau, E. J., Computer challenge corner: Problem 477: A brute force program . J. Rec. Math. (1976).

\bigskip
\noindent\small{Source: \url{https://www.erdosproblems.com/307}}
\end{document}
